\section{RizzoLang Syntax and Semantics}

For our implementation of Rizzo, we have designed an ML-like syntax with let bindings, function definitions, and pattern matching. The syntax is shaped by the semantics of Rizzo and the AST structure required for implementing the \citetitle{ullrich_counting_2021} reference counting semantics.

\citeauthor*{rizzo_modal_2025} did not include a concrete syntax for Rizzo in their paper, but \citetitle{ullrich_counting_2021} included syntax for both the $\lambda_{\text{pure}}$ and $\lambda_{\text{RC}}$ intermediate representations. We have leaned on the structure of the $\lambda_{\text{pure}}$ IR for our syntax design, as it was based on the A-normal form from \citetitle{flanagan_essence_1993}, which is a well-known intermediate representation for functional languages. We chose an ML-like syntax because it is familiar to functional programmers and aligns naturally with the constructs from \citetitle{ullrich_counting_2021}.

In this section we first present the concrete syntax of Rizzo, then describe the type system with a focus on the modal types \texttt{Later} and \texttt{Delayed}. We then describe the operational semantics; particularly the eager evaluation strategy and the \texttt{advance}/\texttt{update} mechanism for reactive signals. Finally, we describe the Rizzo heap structure and how it supports efficient reference counting for immutable data.

\subsection{Syntax}

We have extended the base ML-like syntax with constructs specific to Rizzo, such as the signal constructor (\texttt{::}) and the later/delay constructors. The syntax supports top-level definitions, local let bindings, lambda expressions, pattern matching, and type annotations.

\subsubsection{Top-level Declarations}

Programs in Rizzo consist of a sequence of top-level declarations. There are two forms: value bindings and function definitions.

\begin{verbatim}
let x = e                  (* value binding *)
fun f x1 x2 ... xn = e     (* function definition *)
\end{verbatim}

The function definition \texttt{fun f x1 ... xn = e} is syntactic sugar for \texttt{let f = fun x1 xn -> e}. Functions marked with the \texttt{@effectful} annotation are allowed to perform side effects such as registering output signals.

\subsubsection{Expressions}

The core expression forms are:

\begin{verbatim}
e ::= x                           (* variable *)
    | c                           (* constant: (), true, false, 42, "hello" *)
    | fun x -> e                  (* lambda abstraction *)
    | e1 e2                       (* application *)
    | C (e1, ..., en)             (* constructor application *)
    | let x = e1 in e2            (* local binding *)
    | (e1, e2)                    (* tuple *)
    | e1 :: e2                    (* signal/list constructor *)
    | e1 |> e2                    (* pipe operator *)
    | match e with | p1 -> e1 ... (* pattern matching *)
    | if e1 then e2 else e3       (* conditional *)
    | (e : t)                     (* type annotation *)
\end{verbatim}

The pipe operator \texttt{|>} is left-associative and desugars to application, for which \texttt{a |> f} becomes \texttt{f a}. This is particularly useful for chaining signal transformations and for applying delayed functions to later values.

The signal constructor \texttt{::} is right-associative and constructs a signal from a head value and a tail computation. For example, \texttt{5 :: never} creates a constant signal with value 5.

\subsubsection{Patterns}

For patterns we have chosen to keep it simple, by not allowing recursive patterns, but we have included the necessary forms to support destructuring of tuples, signals, and constructors:

\begin{verbatim}
p ::= _                           (* wildcard *)
    | x                           (* variable binding *)
    | c                           (* constant pattern *)
    | (p1, p2)                    (* tuple pattern *)
    | p1 :: p2                    (* signal/cons pattern *)
    | c :: s                      (* string pattern *)
    | C (p1, ..., pn)             (* constructor pattern *)
\end{verbatim}

Like other head and tail structured datatypes, we have chosen to use \texttt{x :: xs} for Signals, to bind the head of the signal to \texttt{x} and the tail (a \texttt{Later} computation) to \texttt{xs}. We also added it for strings, to allow the head character to be bound separately from the tail string. This is useful for combinators that need to process strings character by character.

% \subsubsection{Example: Signal Combinators}

% The following example demonstrates the syntax through common signal combinators:

% \begin{lstlisting}
% (* Create a constant signal *)
% fun const x : 'a -> Signal 'a = x :: never

% (* Create a signal from a Later computation *)
% fun mk_sig later =
%     (fun a -> a :: mk_sig later) |> later

% (* Map a function over a signal *)
% fun map f s =
%     match s with
%     | x :: xs -> f x :: (map f |> xs)

% (* Sample the current value of ys whenever xs updates *)
% fun sample xs ys = map (fun x -> (x, head ys)) xs
% \end{lstlisting}

% The \texttt{map} function demonstrates several key syntactic features: pattern matching on signals with \texttt{x :: xs}, the pipe operator for delayed application \texttt{(map f |> xs)}, and recursive signal construction.


\subsection{Types}\label{section:rizzolang:types}

Rizzo's type system includes primitive types, compound types, and the \texttt{Later} and \texttt{Delayed} modal types that are central to functional reactive programming. In this section we describe the types available in our implementation and acknowledge the current limitations.

% For our implementation, we have chosen to keep the type system simple by not including user-defined algebraic data types. This means that programmers cannot define their own sum types or record types, and must rely on tuples and signals for data structures. This limitation simplifies the implementation of reference counting, as we do not need to handle arbitrary recursive types, but it does restrict expressiveness. Extending the type system with user-defined types is a potential area for future work.

\subsubsection{Primitive Types and compound types}\label{section:rizzolang:types}
The language includes a small set of primitive types: \texttt{Int}, \texttt{Bool}, \texttt{String}, and the unit type \texttt{Unit}. These are simple by themselves, but can be combined using tuples and function types to create more powerful abstractions.

\begin{verbatim}
t ::= Int | Bool | String | Unit   (* primitive types  *)
    | TypeName                     (* named types *)
    | t1 * t2                      (* tuple types      *)
    | t1 -> t2                     (* function types   *)
    | Later t                      (* later type       *)
    | Delayed t                    (* delay type       *)
    | Signal t                     (* signal type      *)
    | t1 (t2, ..., tn)             (* type application *) 
\end{verbatim}

\subsubsection{Modal Types: Later and Delayed}

% reference back to \ref{section:frp:rizzo}

The modal types \texttt{Later} and \texttt{Delayed} are what distinguish Rizzo from ordinary functional languages. As introduced in Section~\ref{section:frp:rizzo}, these types encode temporal properties of values.
\citeauthor*{rizzo_modal_2025} introduces a set of constructors for these modal types, for which the constructors for \texttt{Later} values are:

\begin{itemize}
    \item $\mathsf{never} : \texttt{Later}\,A$ is the delayed computation that never produces any value.
    % (* a computation that never produces a value *)
    \item $\mathsf{wait}\;c : \texttt{Later}\,A$ waits for an input on channel $c : \texttt{Chan}\,A$ to produce a value.
    % (* waits for input on channel \texttt{c : Chan a} *)
    \item $\mathsf{watch}\;s : \texttt{Later}\,A$ watches a partial signal $s : \texttt{Signal}\,(\texttt{Option}\,a)$ and will produce a value when the head of $s$ is updated to a value of the form $\mathsf{Some}(x)$ for some $x: A$. 
    % \texttt{watch s~:~Later a} (* watches a partial signal $\texttt{Signal }(\texttt{Option }a)$ *)
    \item $\mathsf{tail}\;s : \texttt{Later}\,(\texttt{Signal}\,A)$ builds a delayed computation to access the future value of a signal $s : \texttt{Signal}\,A$.
    % \texttt{tail s : Later (Signal a)}   (* gets the tail of signal *)
    \item $\mathsf{laterapp} \; f \; x : \texttt{Later}\,B$ is the applicative action on \texttt{Later} values. Builds a delayed computation to apply $f : \texttt{Delayed}\,(A \to B)$ to $x : \texttt{Later}\,A$ in the future time step when $x$ produces its value.
    % \texttt{laterapp f x}  (* applies \texttt{f~:~Delayed ($a \to b$)} to \texttt{x~:~Later $a$} whenever \texttt{x} becomes available *)
    \item $\mathsf{sync}\; x \; y : \texttt{Later}\,(\texttt{Sync}\,(A,B))$ is the primitive for handling the asynchronous nature of Rizzo. This constructor builds a delayed computation that will produce a value when either $x : \texttt{Later}\,A$ or $y : \texttt{Later}\,B$, or both produce a value. This either-or behaviour is captured by the $\texttt{Sync}\,(A,B)$ type with its constructors $\mathsf{Left}(x')$, $\mathsf{Right}(y')$, and $\mathsf{Both}(x',y')$.
    % \texttt{sync l1 l2 : Later (Sync a b)}   (* synchronizes two later's, producing a value when either or both tick *)
\end{itemize}

and the constructors for the \texttt{Delayed} type are:
\begin{itemize}
    \item $\mathsf{delay}\;e : \texttt{Delayed}\,A$ delays the expression $e$ to a future time step.
    % \texttt{delay e : Delayed a}   (* creates a delayed computation from expression \texttt{e} *)
    \item $\mathsf{ostar} \; f \; x : \texttt{Delayed}\,B$ builds a delayed computation that applies $f: \texttt{Delayed}\,(A \to B)$ to $x : \texttt{Delayed}\,A$ in any future time step. 
    % \texttt{ostar f x}  (* applies \texttt{f~:~Delayed ($a \to b$)} to a delayed value \texttt{x~:~Delayed a} \todo{when it is advanced in a future time step} *)
\end{itemize}


% \paragraph{The Delayed Type.} A value of type \texttt{Delayed a} represents a computation that will be available at \emph{all} future time steps. This universally quantified nature makes \texttt{Delayed} values safe to use across time boundaries. The main constructor is \texttt{delay e : Delayed a}, which creates a delayed computation from an expression \texttt{e}. The expression \texttt{e} is not evaluated immediately, but is instead, suspended until the delayed value is applied to a later value.

% \subsubsection{The Sync Type}

% In \citetitle{rizzo_modal_2025} \cite{rizzo_modal_2025} the \texttt{sum} type is needed to represent the result of synchronizing two \texttt{Later} computations with the \texttt{sync} signal combinator. To simplify the implementation, we have chosen to use a custom \texttt{Sync} type instead, which is a three-way sum indicating which computations produced values when synchronizing two \texttt{Later} computations with \texttt{sync}:

% \begin{verbatim}
% type Sync a b = Left a | Right b | Both (a * b)
% \end{verbatim}

% In practice, the \texttt{sync} combinator is used to combine two \texttt{Later} computations, and the resulting \texttt{Sync} value indicates whether the first, second, or both signals have ticked in the current time step. This allows combinators that use \texttt{sync} to react differently depending on which inputs updated. For example, the \texttt{zip} combinator uses \texttt{sync} to combine two signals:

% \begin{verbatim}
% fun zip xs ys : Signal 'a -> Signal 'b -> Signal ('a * 'b) =
%     let cont = fun mysync ->
%         match mysync with
%         | Left (xs_)      -> zip xs_ ys
%         | Right (ys_)     -> zip xs ys_
%         | Both (xs_, ys_) -> zip xs_ ys_
%     in
%     (head xs, head ys) :: (cont |> sync (tail xs) (tail ys))
% \end{verbatim}

\subsection{Semantics}\label{section:rizzolang:semantics}
% In order for signals to be reachable during update, they are all stored in a signal-only heap. The heap is structured as a list because the ordering is important for making sure signals are updated in the correct order.

% The reactive semantics describe how signals are updated when time is progressed. Described at a high level, the procedure will mark all signals as old, by moving them to the \texttt{Earlier Heap}. We then iterate through all signals in the heap. 
% For each signal it checks if the signal was dependent, either directly or indirectly, on the channel that produced a value and if so the signal is updated otherwise it is skipped. 
% A signal is then updated by evaluating the delayed computation in the tail, resulting in a new signal which is copied back into the old signal. When done, we move the original signal back to the \texttt{Now Heap} and continue with the next signal.


\todo{This section should prepare the reader for why we have structured the heap the way we have and why we have chosen to add a signal constructor to the AST.
We do this by first listing all the semantics we need to adhere to, and how we have done so (this also forms the base of why we have the memory structure we have chosen).}

Our implementation depends on adhering to the reactive semantics of Rizzo as defined in \citetitle{rizzo_modal_2025}~\cite{rizzo_modal_2025}, while also implementing the reference counting semantics from \citetitle{ullrich_counting_2021}~\cite{ullrich_counting_2021}. In this section we describe how we have structured our implementation to meet these requirements, and how the design of the heap and AST supports these semantics.

\subsubsection{Semantics: Rizzo}

In Rizzo the semantics are split into two layers: the standard evaluation semantics for expressions, and the reactive semantics that describe how Rizzo programs react when time is stepped forward. 
The evaluation semantics follow a standard eager evaluation strategy. The only the exception is $\text{delay} \; e$ which lazily evaluates its the argument $e$.
Since the evaluation semantics do not enforce further restrictions, the focus of this section will be the the reactive semantics of Rizzo.

The reactive semantics consists of three parts: 
The \textit{advance semantics} which describe how \texttt{Later} values are advanced to produce a value for the new time step; 
the \textit{update semantics} which, using the advance semantics, describe how each signal is updated in order to bring it into the current time step;
and \text{step semantics} which describe a complete move of all signals from an earlier time step to the current one. For the purpose of this paper, the step semantics is a while-loop wrapped around the update semantics.

%Goal: describe the semantics of Rizzo primitives: Later- and Delayed constructors. Advance semantics?
\todo{Probably not necessary to mention advance a whole lot more}

%Goal: Goal is to explain the update semantics. Explain the now vs. earlier heap.
\todo{To update signals, we iterate all signals in the topological order / the order of data dependency}. 

\todo{To move a signal from the earlier heap to the now heap, we check if its tail was dependent on the channel that produced a value. If it did not then it is moved unchanged, but we set the updated flag to false (we should explain where this flag belongs). If it did depend on the channel, then we \textit{advance} its tail to obtain a new signal. The data of this signal is copied over the original signal. Perform these steps for every signal!}

%Goal: Here comes a digest, what does all of the above mean? 
In Rizzo the reactive semantics require that:
\begin{itemize}
    \item Any new signal must only depend on other signals in the \texttt{Now heap}.
    \item When an input is produced on a channel, all signals with a tail that depends on the clock of that channel must to be updated.
    \item Signals must be updated based on a topological order of their dependencies. If signal B depends on signal A, signal A must appear before signal B in the heap and thus be updated first.
    \item Evaluating the tail of a signal produces a new signal, which must be copied back into the original signal. %This means that the original signal is updated in place, and any references to it remain valid.
    \item If no reference exists to the intermediate signal produced by evaluating the tail, then the intermediate signal must be deallocated to prevent duplicate work.
    \item New signals must be inserted into the \texttt{Now heap}, to prevent it from being evaluated in the current time step.
    \item 
\end{itemize}

%Goal: Okay, so now we know what is required of the heap structure. How do we represent it? -> Linked list, advantages thereof.

Based on these requirements, we have chosen to structure our heap as a linked list. Using a linked list has the advantage of:
% This gives us a simple way to maintain the topological order of signals, it allows for efficient insertion and removal of signals, 
%  as we can ensure that when a new signal is created it is added after the split between the \texttt{Now} and \texttt{Earlier} heaps. \dots

\begin{itemize}
    \item Representing the \texttt{Now} and \texttt{Earlier} heaps in one data structure, and expressing a split between them using a pointer (cursor).
    \item Moving signals from the \texttt{Now} heap to the \texttt{Earlier} heap in O(1) time by resetting the cursor to the start of the heap. Likewise moving signals back to the \texttt{Now} heap by moving the cursor forward.
    \item O(1) insertion that maintaining the topological order of signals, by adding them in the end of the \texttt{Now} heap, we ensure that any signal can only depend on signals that were created before it.
    \item O(1) removal of signals when their reference count drops to zero.
    \item Stable references to signals, as insertions and removals do not move other signals in memory, so pointers to signals remain valid throughout their lifetime.
\end{itemize}

\subsubsection{How does reference counting fit}

\todo{Where do we reference count in the reactive semantics of Rizzo? The intermediate signal created inside update is decremented after the data has been written over \textemdash\ to allow immediate erasure \dots}

\subsubsection{Added Constructors for Rizzo}
%Operational semantics: signals are different from normal constructors, they should go in the signal heap ... That is the goal below?

In the counting immutable beans semantics, it says that a constructor call should return a location to the object that has been created. For Rizzo, we have extended this semantic for signals, by also having the signal constructor add the signal to the signal heap.

\todo{We should mention that all constructors are heap allocated. We could show example of general heap to show how the signal heap is just a subset.}



% For RizzoLang we need to adhere to the \texttt{advance} and \texttt{update} semantics of Rizzo as defined in \citetitle{rizzo_modal_2025} \cite{rizzo_modal_2025}, but we also need to define the semantics of reference counting for our implementation. In this section we first describe the evaluation strategy of Rizzo, then the reactive semantics for signals, and finally how reference counting interacts with the heap structure.

% % Rizzo has two layers of semantics: the standard evaluation semantics for expressions, and the reactive semantics that govern how signals are updated over time. Understanding both is essential for implementing an efficient runtime.

% \subsubsection{Eager Evaluation}

% Rizzo uses eager (strict) evaluation for all expressions. Function arguments are evaluated before the function body executes, and let bindings evaluate their right-hand side before the body. This is in contrast to lazy languages like Haskell where evaluation is deferred until values are needed.

% However, there is an important exception: \texttt{Delayed} and \texttt{Later} computations are not evaluated immediately. When we write \texttt{delay e}, the expression \texttt{e} is not evaluated---instead, a thunk (suspended computation) is created. Similarly, when we write \texttt{f |> later}, the application of \texttt{f} is deferred until the \texttt{later} value becomes available.

% This combination of eager evaluation with explicit delayed computations gives programmers precise control over when computations occur. Signal combinators use \texttt{delay} to ensure that recursive calls happen only when time advances, preventing infinite loops.

% \subsubsection{Reactive Semantics: Time Steps and Updates}

% The reactive semantics describe how signals evolve over time. Time in Rizzo is discrete: it advances in steps, with each step triggered by input arriving on a channel. Between time steps, signals maintain their current values.

% When a channel produces a value, the runtime must update all signals that depend (directly or indirectly) on that channel. This is done through two key operations: \texttt{ticked} and \texttt{advance}.

% \paragraph{The Ticked Predicate.} Given a \texttt{Later} computation and a channel that just produced a value, the \texttt{ticked} predicate determines whether the computation should produce a value in this time step:

% \begin{itemize}
%     \item $\texttt{ticked}(\texttt{never}, c, v) = \texttt{false}$ --- never produces a value
%     \item $\texttt{ticked}(\texttt{wait } k, c, v) = (k = c)$ --- ticks when the waited channel fires
%     \item $\texttt{ticked}(\texttt{tail } s, c, v) = \texttt{updated}(s)$ --- ticks when the signal was updated
%     \item $\texttt{ticked}(\texttt{sync } l_1 \; l_2, c, v) = \texttt{ticked}(l_1, c, v) \lor \texttt{ticked}(l_2, c, v)$
%     \item $\texttt{ticked}(f \triangleright l, c, v) = \texttt{ticked}(l, c, v)$ --- follows the later argument
% \end{itemize}

% \paragraph{The Advance Operation.} When a \texttt{Later} computation has ticked, the \texttt{advance} operation extracts the value it produces:

% \begin{itemize}
%     \item $\texttt{advance}(\texttt{wait } k, c, v) = v$ --- returns the channel's value
%     \item $\texttt{advance}(\texttt{tail } s, c, v) = s$ --- returns the (now updated) signal
%     \item $\texttt{advance}(f \triangleright l, c, v) = f(\texttt{advance}(l, c, v))$ --- applies the delayed function
%     \item $\texttt{advance}(\texttt{sync } l_1 \; l_2, c, v)$ returns \texttt{Left}, \texttt{Right}, or \texttt{Both} depending on which computations ticked
% \end{itemize}

% \paragraph{The Heap Update Procedure.} The runtime maintains all signals in a heap (see Section~\ref{section:rizzolang:heap}). When a channel fires, the update procedure iterates through the heap:

% \begin{enumerate}
%     \item For each signal $s$ with tail $t$:
%     \item If $\texttt{ticked}(t, c, v)$:
%     \begin{enumerate}
%         \item Compute $s' = \texttt{advance}(t, c, v)$ --- this produces a new signal
%         \item Copy $\texttt{head}(s')$ and $\texttt{tail}(s')$ back into $s$
%         \item Mark $s$ as updated
%         \item The intermediate signal $s'$ can now be deallocated
%     \end{enumerate}
%     \item Otherwise, mark $s$ as not updated
% \end{enumerate}

% The copying step is crucial: it ensures that references to $s$ remain valid while the signal's contents are updated. This is the ``mutable reference'' semantics that distinguishes Rizzo from other FRP languages.
% % this will become important as to why we cant reuse the intermediate signal

% \subsubsection{Update Ordering and Cascading}

% The order in which signals are updated matters. Consider a signal $s_2$ whose tail depends on another signal $s_1$. If $s_1$ updates, then $s_2$ may also need to update (if its tail uses \texttt{tail} $s_1$). For correct semantics, $s_1$ must be updated before $s_2$.

% The heap maintains signals in creation order, which ensures that dependencies are respected: a signal can only depend on signals created before it. When a new signal is created during an update (e.g., by advancing a \texttt{Later} computation), it is inserted into the now heap, ensuring the signal will not be updated until the next time step. This allows for cascading updates without risking infinite loops within a single time step.

% \subsubsection{Reference Counting Semantics}

% The \citetitle{ullrich_counting_2021} paper defines a reference-counted intermediate representation ($\lambda_{\text{RC}}$) with explicit \texttt{inc}, \texttt{dec}, \texttt{reset}, and \texttt{reuse} operations. Our compiler transforms Rizzo programs into this representation before generating C code.

% The key semantic rules are:

% \begin{itemize}
%     \item Every variable starts with a reference count of 1 when bound
%     \item Passing a variable to an \emph{owned} parameter consumes one reference
%     \item Passing a variable to a \emph{borrowed} parameter does not consume a reference
%     \item \texttt{inc x} increments the reference count of \texttt{x}
%     \item \texttt{dec x} decrements the reference count; if it reaches zero, the object is freed
%     \item \texttt{reset x} prepares \texttt{x} for potential reuse if its count is 1
%     \item \texttt{reuse x in C(...)} reuses the memory of \texttt{x} for a new constructor
% \end{itemize}

% These operations are inserted automatically by the compiler based on ownership analysis. The details of this transformation are covered in Section~\ref{section:refcount}.

\section{The Rizzo Heap}\label{section:rizzolang:heap}
The reactive semantics of Rizzo require that all signals be reachable during the update phase. To support this, signals are stored in a dedicated heap structure. In this section we describe the design of this heap and how it interacts with reference counting to enable efficient memory management.

\subsection{Object Representation}

All heap-allocated objects in Rizzo share a common header structure:

\begin{verbatim}
struct rz_header {
    uint16_t num_fields;    /* number of fields in the object */
    uint8_t  tag;           /* constructor tag for pattern matching */
    uint8_t  obj_type;      /* object kind: regular, string, signal, etc. */
    int32_t  refcount;      /* reference count */
};
\end{verbatim}

The \texttt{num\_fields} and \texttt{tag} fields support pattern matching and recursive deallocation. The \texttt{obj\_type} field distinguishes between regular objects, strings, signals, and partial applications---each requiring different handling during deallocation. The \texttt{refcount} field tracks how many references point to this object.

Following the header, objects store their fields as an array of boxed values. A boxed value (\texttt{rz\_box\_t}) can hold either an immediate integer, a string literal pointer, or a pointer to a heap object. This uniform representation simplifies the generated code at the cost of some memory overhead.

\subsection{Signal Structure}

Signals have additional fields beyond the standard header:

\begin{verbatim}
struct rz_signal {
    rz_header_t header;
    rz_box_t head;      /* current value */
    rz_box_t tail;      /* Later computation for future values */
    rz_box_t updated;   /* flag: was this signal updated this tick? */
    rz_box_t prev;      /* previous signal in heap (doubly-linked list) */
    rz_box_t next;      /* next signal in heap */
};
\end{verbatim}

The \texttt{head} and \texttt{tail} fields correspond to the signal definition from Section~\ref{section:frp:rizzo}. The \texttt{updated} flag is set during the heap update procedure and is used by combinators that need to know whether a signal changed (e.g., \texttt{watch}).

The \texttt{prev} and \texttt{next} pointers form a doubly-linked list, which is the heap structure itself.

\subsection{The Heap as a Linked List}

The signal heap is implemented as a doubly-linked list rather than a contiguous array. This design choice is motivated by several requirements:

\begin{enumerate}
    \item \textbf{Ordered iteration}: Signals must be updated in the order they were created to respect dependencies. For example, if signal $s_2$ depends on $s_1$, then $s_1$ must be updated before $s_2$. A linked list naturally preserves insertion order, while an array would require additional bookkeeping to maintain order during insertions and deletions.
    % Signals must be updated in creation order to respect dependencies. A linked list naturally preserves insertion order.

    \item \textbf{Efficient insertion}: New signals created during an update must be inserted at the current position in the iteration \todo{in the now heap?}. With a linked list, this is O(1).

    \item \textbf{Efficient removal}: When a signal's reference count drops to zero, it must be removed from the heap. Doubly-linked lists support O(1) removal.

    \item \textbf{Stable references}: Unlike arrays, linked lists do not move elements when inserting or removing, so pointers to signals remain valid.
\end{enumerate}

The runtime maintains two pointers into the heap:
\begin{itemize}
    \item \texttt{heap\_base}: Points to the first signal in the heap
    \item \texttt{heap\_cursor}: Points to the signal currently being processed during update
\end{itemize}

\subsection{Signal Lifecycle}

A signal goes through the following lifecycle:

\paragraph{Creation.} When a signal is constructed (via \texttt{::}), it is allocated and inserted into the heap just before the current cursor position. This ensures that newly created signals are not considered for update in the current time step. The signal starts with a reference count of 1.

\paragraph{Update.} During the heap update phase, if a signal's tail has ticked, the runtime:
\begin{enumerate}
    \item Advances the tail to produce a new (intermediate) signal
    \item Copies the head and tail from the intermediate signal to the original
    \item Decrements the reference count of the old head and tail
    \item Increments the reference count of the new head and tail
    \item Decrements the reference count of the intermediate signal
\end{enumerate}

If the intermediate signal's reference count reaches zero (which is typical), it is immediately deallocated and removed from the heap.

\paragraph{Deallocation.} When a signal's reference count reaches zero:
\begin{enumerate}
    \item The reference counts of its head and tail are decremented (potentially triggering further deallocations)
    \item The signal is removed from the linked list by updating its neighbors' pointers
    \item The signal's memory is freed
\end{enumerate}

\subsection{Why Reference Counting Matters for the Heap}

The heap update procedure creates intermediate signals that are typically used only to copy values back to the original signal. Without proper memory management, these intermediate signals would accumulate in the heap, causing two problems:

\begin{enumerate}
    \item \textbf{Memory leaks}: Intermediate signals consume memory indefinitely.

    \item \textbf{Redundant work}: During the next update, the runtime would iterate over and potentially update these orphaned signals, wasting CPU cycles and potentially causing incorrect behavior.
\end{enumerate}

Reference counting solves both problems by ensuring that intermediate signals are deallocated immediately when they are no longer needed. The key insight is that an intermediate signal typically has a reference count of 1 (held by the local variable in the update procedure). When that reference is released, the signal is freed and removed from the heap before the next iteration.

This is why the \citetitle{ullrich_counting_2021} approach is particularly well-suited to Rizzo: the combination of precise reference counting with the \texttt{reset}/\texttt{reuse} optimization allows the runtime to efficiently recycle signal memory, minimizing both allocation overhead and heap size.

\subsection{Memory Layout Summary}

To summarize, the Rizzo runtime manages memory at two levels:

\begin{enumerate}
    \item \textbf{Object level}: All heap objects (including signals) use reference counting for lifetime management. The compiler inserts \texttt{inc}/\texttt{dec} operations, and the runtime performs deallocation when counts reach zero.

    \item \textbf{Heap level}: Signals are additionally linked into the signal heap for iteration during updates. The heap structure is maintained automatically as signals are created and destroyed.
\end{enumerate}

This two-level design separates concerns: reference counting handles object lifetimes, while the heap structure supports the reactive semantics. The detailed implementation of reference count insertion is covered in Section~\ref{section:refcount}.

